%% 06-the-decision-engine.tex -- generated by tools/build_from_docx.py; edit freely

\hypertarget{the-decision-engine}{%
\section{The Decision Engine}\label{the-decision-engine}}

The Decision Engine is the capability through which an agent decides
what shall change. It consumes insight, objectives, policy and current
state, and emits intent: a desired outcome, with the scope over which it
applies, the constraints under which it holds and the window through
which it remains valid. It is the capability that carries authority, and
it is therefore the capability the framework constrains most tightly.

Two boundaries define it. The Decision Engine emits intent and never a
command: translation of intent into device-specific configuration
belongs to the Actuator, and a Decision Engine emitting
controller-native payloads would have absorbed the Actuator's guardrails
into a capability that does not implement them. An intent contract whose
schema carries controller-native command fields is rejected at
admission.

The second boundary is the one on which the framework's safety rests,
and it admits no exception: a Decision Engine cannot approve its own
decision. Every intent affecting a shared or contended resource is
mediated by the conflict fabric before it can reach an Actuator.
Coordination is never internal to a decision. This is what makes
federation across domains safe in the absence of any global reasoning
entity, and it holds for every Decision Engine without exception,
including experimental ones.

\hypertarget{topology}{%
\subsection{Topology}\label{topology}}

Decision Engines are per domain. A cross-domain agent is a federation of
per-domain closed loops, each deciding within its own authority
envelope, bounded by the delegation its domain received. There is no
runtime global Decision Engine, and the framework does not admit one.

End-to-end coherence across domains comes from two mechanisms and not
from a central decider. When a global agent is instantiated, its
end-to-end objective is resolved once into per-domain portions, which
establishes what each domain is to achieve. At runtime the conflict
fabric arbitrates between decisions that contend, whether over a shared
control parameter or between objectives that pull against each other.
The first mechanism aligns objectives before anything runs; the second
resolves contention while it runs.

Decisions may be exchanged between domain agents. Such an exchange is a
proposal: a decision received from a peer binds nothing, confers no
authority over the receiving agent, and leaves the receiving agent
answerable for whatever it applies within its own scope. Where a binding
decision across domains is required, it is taken by the overlooking
agent at the cross-domain management level. Authority descends; it does
not pass sideways.

\hypertarget{capability-declarations-2}{%
\subsection{Capability Declarations}\label{capability-declarations-2}}

At \texttt{full} the Decision Engine decides within its declared action
space. At \texttt{partial} it is restricted to advisory or bounded
decisions, and the restriction on its decision space shall be
enumerated; an unenumerated partial declaration is rejected at
admission. At \texttt{delegated} the decision is taken by another agent
or service, which is named. At \texttt{none} the agent holds no decision
capability and is a sensor, an analytics function or an actuator alone.

The deployment tier follows the decision latency the loop requires,
bounded by the effective latency class computed at admission. A radio
access Decision Engine commonly deploys as an application on the
near-real-time controller, which has a decisive consequence for how
conflict mediation is realized; a transport Decision Engine deploys
regionally or centrally against the transport controller.

\hypertarget{the-egress-gate}{%
\subsection{The Egress Gate}\label{the-egress-gate}}

\hypertarget{mediation-as-a-property-of-the-runtime}{%
\subsubsection{Mediation as a property of the runtime}\label{mediation-as-a-property-of-the-runtime}}

The Decision Engine runtime does not publish intent to an Actuator.
Intent leaves the capability through the conflict fabric binding, which
returns one of four verdicts: allow, modify, defer or deny. The binding
is established at activation and is a precondition of the agent being
granted decision emission at all; an agent whose conflict binding is
unresolved is never granted emission.

Not every intent requires arbitration. Where the parameters an intent
touches lie in an action space held exclusively by one agent, and no key
performance indicator they affect is registered as another agent's
objective, there is no contention to resolve, and the fabric returns
allow without an arbitration cycle. Exclusive control of a parameter is
not by itself freedom from conflict: two agents acting on distinct
parameters conflict in their intents when the indicators those
parameters move are objectives the agents pursue against each other. The
determination is made by the platform from the action spaces and
objectives declared in the intent contracts of the admitted agents,
never by the emitting agent about its own intent; an agent permitted to
judge its own decision uncontested would hold self-approval under
another name. Exclusivity is therefore a property the platform computes
and the agent cannot assert. Where it holds, the mediation step is
structurally absent rather than skipped, which is what keeps the fast
path viable for the tightest latency classes.

The Decision Engine's work is unchanged by this distinction. It forms
intents, accumulates them against its objectives and hands them to
actuation whether or not an arbitration cycle intervened; mediation
governs egress, not the formation of decisions.

This is what converts the prohibition on self-approval from a rule into
a property. A rule can be disregarded by an agent that does not
implement it; a physical inability to reach the Actuator without a
verdict cannot. Mediation is an egress property of the platform runtime
rather than a step the agent is trusted to perform, and no
misconfiguration, no experimental status and no partner implementation
choice can bypass it.

The contract is identical across domains; only its realization differs.
Where a Decision Engine is deployed as an application on the
near-real-time controller, the control actions it emits are observed
natively by the conflict management function of that controller. Where a
domain has no such controller, as in transport, mediation is an explicit
synchronous call to the conflict management function local to the domain
management function before egress. The descriptor expresses the choice
through the conflict resolution policy already defined for the framework
--- priority, arbitration, negotiation or hierarchical, with priority
rules --- and requires no additional field.

\hypertarget{bounded-modification}{%
\subsubsection{Bounded modification}\label{bounded-modification}}

Where the fabric returns modify, the modified intent is re-checked
against the agent's own declared action space and constraints before it
proceeds. A modification falling outside those bounds is treated as a
denial, and the agent's declared behavior for unresolved conflict
executes.

Without this check, arbitration would become a route around the bounds
an agent was admitted under. The fabric could hand an agent an intent it
would never have been permitted to form, and the agent would apply it on
the fabric's authority. The re-check keeps the agent's admitted envelope
binding on everything the agent emits, whatever the source of the
modification.

\hypertarget{stability-controls}{%
\subsubsection{Stability controls}\label{stability-controls}}

Declared stability controls --- cooldown, hysteresis, maximum rate of
change and staging --- are enforced by the platform at the egress gate
and again on the actuation path. A publisher declaring a thirty-second
cooldown whose logic emits every second has its excess intents
suppressed at egress, and a conformance fault is recorded against the
agent.

The general principle is the one the framework applies throughout: a
declared bound is a guarantee only where something other than the
declaring party enforces it. A cooldown the agent enforces on itself is
a description of the agent's intentions, not a property of the system.

\hypertarget{evidence-gating}{%
\subsection{Evidence Gating}\label{evidence-gating}}

The Decision Engine runtime evaluates the input quality of the insight
it consumed before the decision logic runs. Where the degraded-input
flag is set, or coverage falls below the minimum the agent declared, the
declared degraded-evidence behavior executes and the contributed
decision logic is not invoked at all for that cycle.

The ordering is the substance of the rule. If evaluation of input
quality were left inside contributed logic, a publisher could ignore the
field, and the chain that begins with coverage computed at the
Monitoring System would terminate one step short of the capability that
acts. Gating before the logic runs makes the check unskippable, and
completes the propagation of evidence quality across the whole loop: the
Monitoring System computes coverage against a resolved scope, the
Analytics Engine derives input quality from the envelopes consumed, and
the Decision Engine acts on it structurally rather than by the model's
cooperation.

\hypertarget{decision-rationale}{%
\subsection{Decision Rationale}\label{decision-rationale}}

Every emitted intent carries a structured rationale: the identifiers of
the insight envelopes consumed, the constraints applied, the objectives
targeted and a reason code. Free text alone is not conformant, and an
intent without a rationale is rejected at the publication boundary.

The framework requires traceable provenance of every agent and model
interpretability of those holding the highest authority. Provenance is
achievable universally: every agent supplies the inputs it consumed, the
constraints it applied and a reason code, because those are properties
of the decision's context rather than of the model's internals.
Interpretability is not universally achievable, and requiring it of
every agent would exclude most contributed algorithms and with them most
of the work the platform exists to host.

The descriptor therefore carries interpretability as a declared property
of the decision logic, stating whether the agent can produce, for an
individual decision, an account of which inputs determined the outcome
and in what direction. An agent that declares it supplies that account
alongside the rationale on request. An agent that does not declare it
remains admissible, and is not eligible for
closed\_\allowbreak{}loop\_\allowbreak{}full authority; its highest
attainable mode is auto\_\allowbreak{}execute\_\allowbreak{}guarded,
where an external guard stands between its intent and the network.
Authority over the network without an account of why a decision was
taken is the combination the framework does not permit; either property
alone is admissible at the authority that matches it.

What rationale makes possible is threefold: audit after the fact,
arbitration between agents whose intents contend, and the assurance
record on which an agent's continued authority rests.

\hypertarget{actuation-feedback-and-the-re-decision-budget}{%
\subsection{Actuation Feedback and the Re-Decision Budget}\label{actuation-feedback-and-the-re-decision-budget}}

The Actuator publishes typed outcomes --- acknowledged and verified,
acknowledged but unverified, rejected, or unresolved --- and the
Decision Engine consumes them as evidence. An intent that was rejected
or could not be verified is information about the network the agent did
not have when it decided, and re-deciding on that basis is legitimate:
the decision that was taken is no longer available, and the agent
reassesses.

Re-decision is bounded rather than forbidden. The failure policy
declares a maximum number of re-decisions per objective per window, and
the counter is held by the platform runtime rather than by the agent. On
exhaustion the agent stops and raises a lifecycle event to its domain
management function recording that the objective is unachievable,
instead of continuing to decide.

This closes the oscillation in which an agent decides, fails, and
decides again indefinitely, without forbidding the re-evaluation that a
changed network legitimately requires. An objective that cannot be
achieved becomes a matter for the management layer rather than a loop
that runs forever.

\hypertarget{declared-failure-behavior-2}{%
\subsection{Declared Failure Behavior}\label{declared-failure-behavior-2}}

The Decision Engine's failure policy declares the five responses of
Table 6-1.

It declares the response to low confidence, as a safe intent, no action,
or a request for arbitration, together with the confidence threshold
that triggers it. It declares the response to degraded evidence, which
shall be declared independently of the confidence threshold, since the
two express different things and a model confident on degraded evidence
is exactly the case the framework guards against. It declares the
response to unresolved conflict on arbitration timeout, as deferral, a
safe intent, or abandonment. It declares the response to actuation
failure, including the maximum re-decision count per objective per
window. And it declares the response when no usable insight is
available.

A hard invariant governs all five: where a declared failure behavior is
ambiguous or absent at runtime, the Decision Engine emits no intent.
Silence is the safe default, and no action is always a permitted
outcome. An agent that cannot establish what it ought to do does
nothing, and the framework never treats inaction as a failure requiring
a remedy.

\hypertarget{experimental-decision-engines}{%
\subsection{Experimental Decision Engines}\label{experimental-decision-engines}}

A Decision Engine whose logic cannot be verified against the normative
schema is admitted as an extension, at reduced authority: a narrower
control scope, tighter action space bounds and lower arbitration
priority. It uses the identical egress path, the identical conflict
mediation, the identical stability enforcement and the identical
rationale requirements as any other agent.

The freedom applies to the algorithm and never to the guardrails. An
escape hatch exists so that an unsupported algorithm can be tried on the
platform; it does not exist to obtain weaker safety. Containment that
varied with conformance would make the hatch the cheapest route past the
constraints, which is precisely the drift it must not cause.

Sustained conformant operation under evaluation is the route to
promotion into the normative schema and to full authority. The extension
path is therefore a pipeline into the standard rather than a permanent
margin outside it.

\begin{wptable}{6-1}

\begin{longtable}[]{@{}
  >{\raggedright\arraybackslash}p{(\columnwidth - 6\tabcolsep) * \real{0.2875}}
  >{\raggedright\arraybackslash}p{(\columnwidth - 6\tabcolsep) * \real{0.1250}}
  >{\raggedright\arraybackslash}p{(\columnwidth - 6\tabcolsep) * \real{0.1000}}
  >{\raggedright\arraybackslash}p{(\columnwidth - 6\tabcolsep) * \real{0.4874}}@{}}
\caption{Decision Engine failure policy keys}\tabularnewline
\toprule\noalign{}
\begin{minipage}[b]{\linewidth}\raggedright
\textbf{Field in failure\_\allowbreak{}policy.\allowbreak{}DE}
\end{minipage} & \begin{minipage}[b]{\linewidth}\raggedright
\textbf{Type}
\end{minipage} & \begin{minipage}[b]{\linewidth}\raggedright
\textbf{Req.}
\end{minipage} & \begin{minipage}[b]{\linewidth}\raggedright
\textbf{Definition}
\end{minipage} \\
\midrule\noalign{}
\endfirsthead
\toprule\noalign{}
\begin{minipage}[b]{\linewidth}\raggedright
\textbf{Field in failure\_\allowbreak{}policy.\allowbreak{}DE}
\end{minipage} & \begin{minipage}[b]{\linewidth}\raggedright
\textbf{Type}
\end{minipage} & \begin{minipage}[b]{\linewidth}\raggedright
\textbf{Req.}
\end{minipage} & \begin{minipage}[b]{\linewidth}\raggedright
\textbf{Definition}
\end{minipage} \\
\midrule\noalign{}
\endhead
\bottomrule\noalign{}
\endlastfoot
\textbf{low\_\allowbreak{}confidence} & enum & M &
safe\_\allowbreak{}intent, no\_\allowbreak{}action or
request\_\allowbreak{}arbitration, with the confidence threshold that
triggers it. \\
\textbf{degraded\_\allowbreak{}evidence} & enum & M & Behavior when the
degraded-input flag is set. Declared independently of the confidence
threshold. \\
\textbf{conflict\_\allowbreak{}unresolved} & enum & M & defer,
safe\_\allowbreak{}intent or abandon, on arbitration timeout. \\
\textbf{actuation\_\allowbreak{}failure\_\allowbreak{}response} & object
& M & Behavior on rejected, unverified or unresolved outcomes, including
the maximum re-decisions per objective per window. \\
\textbf{insight\_\allowbreak{}unavailable} & enum & M & Behavior when
the AE supplies nothing usable. \\
\end{longtable}

\end{wptable}

\hypertarget{the-intent-contract}{%
\subsection{The Intent Contract}\label{the-intent-contract}}

An intent contract declares what a Decision Engine emits: the objectives
it may pursue, the action space within which its intents fall, the
constraints it holds itself to, the stability controls it declares and
the conflict resolution policy under which it is mediated. Its fields
are given in Table 6-2.

\begin{wptable}{6-2}

\begin{longtable}[]{@{}
  >{\raggedright\arraybackslash}p{(\columnwidth - 6\tabcolsep) * \real{0.2875}}
  >{\raggedright\arraybackslash}p{(\columnwidth - 6\tabcolsep) * \real{0.1250}}
  >{\raggedright\arraybackslash}p{(\columnwidth - 6\tabcolsep) * \real{0.1499}}
  >{\raggedright\arraybackslash}p{(\columnwidth - 6\tabcolsep) * \real{0.4375}}@{}}
\caption{Intent contract fields}\tabularnewline
\toprule\noalign{}
\begin{minipage}[b]{\linewidth}\raggedright
\textbf{Field in intent\_\allowbreak{}contract}
\end{minipage} & \begin{minipage}[b]{\linewidth}\raggedright
\textbf{Type}
\end{minipage} & \begin{minipage}[b]{\linewidth}\raggedright
\textbf{Req.}
\end{minipage} & \begin{minipage}[b]{\linewidth}\raggedright
\textbf{Definition}
\end{minipage} \\
\midrule\noalign{}
\endfirsthead
\toprule\noalign{}
\begin{minipage}[b]{\linewidth}\raggedright
\textbf{Field in intent\_\allowbreak{}contract}
\end{minipage} & \begin{minipage}[b]{\linewidth}\raggedright
\textbf{Type}
\end{minipage} & \begin{minipage}[b]{\linewidth}\raggedright
\textbf{Req.}
\end{minipage} & \begin{minipage}[b]{\linewidth}\raggedright
\textbf{Definition}
\end{minipage} \\
\midrule\noalign{}
\endhead
\bottomrule\noalign{}
\endlastfoot
\textbf{contract\_\allowbreak{}id} & string & M & Identifier of the
intent contract. \\
\textbf{version} & string & M & Version of the contract. \\
\textbf{objectives{[}{]}} & array & M & Objectives the DE may pursue. \\
\textbf{action\_\allowbreak{}space} & object & M & Managed objects and
attributes the DE may seek to change, with permitted ranges.
Controller-native command fields are rejected. \\
\textbf{constraints{[}{]}} & array & M & Constraints the DE holds itself
to, as structured predicates. \\
\textbf{stability\_\allowbreak{}controls} & object & M & Cooldown,
hysteresis, maximum rate of change and staging. Enforced externally. \\
\textbf{conflict\_\allowbreak{}resolution} & enum & M & priority,
arbitration, negotiation or hierarchical, with priority rules. \\
\textbf{decision\_\allowbreak{}space\_\allowbreak{}restriction} & object
& C: mode is partial & Enumerated restriction on the decision space.
Unenumerated partial is rejected. \\
\textbf{validity\_\allowbreak{}window} & duration & M & Window through
which an emitted intent remains valid. \\
\textbf{schema\_\allowbreak{}ref} & string & M & Semantic schema of the
intent payload, with its version. \\
\end{longtable}

\end{wptable}

Table A-1 gives the descriptor location at which each field of Tables
6-1 and 6-2 is declared.

A contract whose schema carries controller-native command fields ---
device configuration payloads, controller operation names, vendor
parameter identifiers --- is rejected at admission, whatever else it
declares. The action space is expressed in terms of managed objects and
their attributes, not in terms of the operations a particular controller
exposes, because an action space stated in a controller's own vocabulary
could not be checked against an agent's admitted bounds by anything
other than that controller.

\hypertarget{the-intent-envelope}{%
\subsection{The Intent Envelope}\label{the-intent-envelope}}

Every intent carries the envelope of Table 6-3, populated by the
Decision Engine runtime. An intent reaching the egress gate without a
complete envelope is rejected there rather than mediated.

\begin{wptable}{6-3}

\begin{longtable}[]{@{}
  >{\raggedright\arraybackslash}p{(\columnwidth - 6\tabcolsep) * \real{0.2875}}
  >{\raggedright\arraybackslash}p{(\columnwidth - 6\tabcolsep) * \real{0.1250}}
  >{\raggedright\arraybackslash}p{(\columnwidth - 6\tabcolsep) * \real{0.1499}}
  >{\raggedright\arraybackslash}p{(\columnwidth - 6\tabcolsep) * \real{0.4375}}@{}}
\caption{Intent envelope fields}\tabularnewline
\toprule\noalign{}
\begin{minipage}[b]{\linewidth}\raggedright
\textbf{Field in envelope}
\end{minipage} & \begin{minipage}[b]{\linewidth}\raggedright
\textbf{Type}
\end{minipage} & \begin{minipage}[b]{\linewidth}\raggedright
\textbf{Req.}
\end{minipage} & \begin{minipage}[b]{\linewidth}\raggedright
\textbf{Definition}
\end{minipage} \\
\midrule\noalign{}
\endfirsthead
\toprule\noalign{}
\begin{minipage}[b]{\linewidth}\raggedright
\textbf{Field in envelope}
\end{minipage} & \begin{minipage}[b]{\linewidth}\raggedright
\textbf{Type}
\end{minipage} & \begin{minipage}[b]{\linewidth}\raggedright
\textbf{Req.}
\end{minipage} & \begin{minipage}[b]{\linewidth}\raggedright
\textbf{Definition}
\end{minipage} \\
\midrule\noalign{}
\endhead
\bottomrule\noalign{}
\endlastfoot
\textbf{intent\_\allowbreak{}id} & uuid & M & Unique identifier of the
intent, assigned by the DE runtime. \\
\textbf{agent\_\allowbreak{}instance\_\allowbreak{}id} & string & M &
Instance that formed the intent. \\
\textbf{contract\_\allowbreak{}id} & string & M & Intent contract and
version under which the intent is emitted. \\
\textbf{objective\_\allowbreak{}ref} & string & M & Objective the intent
pursues. \\
\textbf{scope} & object & M & Managed objects over which the intent
applies, within the agent\textquotesingle s action space. \\
\textbf{desired\_\allowbreak{}outcome} & object & M & The outcome
sought, expressed over managed object attributes. \\
\textbf{constraints\_\allowbreak{}applied{[}{]}} & array & M &
Constraints the DE applied in forming the intent. \\
\textbf{validity\_\allowbreak{}window} & duration & M & Window through
which the intent remains valid. \\
\textbf{urgency} & integer & O & Priority claim ordering the agent's own
intents within its registered priority band (Chapter 10); never raises
the band. \\
\textbf{candidates{[}{]}} & array & O & Alternative outcomes meeting the
same objective, each within the action space; a modify verdict may name
the one the fabric permits (Chapter 10). \\
\textbf{rationale} & object & M & Insight envelope identifiers consumed,
constraints applied, objectives targeted and a reason code. Free text
alone is not conformant. \\
\textbf{input\_\allowbreak{}quality\_\allowbreak{}ref} & object & M &
Input quality of the insight consumed, carried forward from the AE. \\
\textbf{verdict} & enum & M & allow, modify, defer or deny, written by
the runtime after mediation and never by the agent. \\
\textbf{modified\_\allowbreak{}intent} & object & C: verdict is modify &
The intent as modified by the fabric, carried alongside the intent as
formed. \\
\textbf{recheck\_\allowbreak{}outcome} & enum & C: verdict is modify &
Result of the re-check of a modified intent against the
agent\textquotesingle s own bounds. \\
\textbf{redecision\_\allowbreak{}index} & integer & M & Position within
the re-decision budget for this objective and window, held by the
runtime. \\
\textbf{decided\_\allowbreak{}time} & timestamp & M & Time the intent
was formed. \\
\end{longtable}

\end{wptable}

The envelope records what the agent decided, on what evidence, under
which constraints, and what the fabric returned. The verdict field is
written by the runtime after mediation and never by the agent, so that
the record of what was permitted is not the record of what was
requested. Where the verdict was a modification, both the intent as
formed and the intent as modified are carried, together with the outcome
of the re-check against the agent's own bounds.

\hypertarget{requirements-3}{%
\subsection{Requirements}\label{requirements-3}}

The requirements of this section apply to every agent that declares a
Decision Engine in any mode other than none, and to the platform that
admits and runs it.

\begin{requirement}{Req-T23-DE-01}{MUST}{Intent, never command}
\reqfield{Description / rationale}{%
The DE shall emit intent expressed as a desired outcome with scope,
constraints and a validity window. An intent contract whose schema
carries controller-native command fields shall be rejected at admission,
and the action space shall be expressed over managed objects and their
attributes rather than over controller operations.}
\reqfield{Dependency}{%
Req-T23-FW-02}
\reqfield{Traceability (Forward)}{%
\begin{itemize}
\tightlist
\item
  Checked at admission against the registered intent schema.
\item
  Enforced at the egress gate.
\end{itemize}}
\reqfield{UNITY-6G components}{%
\begin{itemize}
\tightlist
\item
  Tuple descriptor
\item
  DE runtime
\item
  Admission (IDMO/DMO)
\end{itemize}}
\reqfield{Linked KPIs}{%
\begin{itemize}
\tightlist
\item
  KPI\#1: Deliver four distributed management entities (MS, AE, DE and
  ACT) deployable in each heterogeneous domain.
\item
  KPI\#42: DT-validated AI algorithms that do not break the
  time-sensitive operation when deployed in 99\% of the outputs.
\end{itemize}}
\end{requirement}

\begin{requirement}{Req-T23-DE-02}{MUST}{No self-approval}
\reqfield{Description / rationale}{%
A DE shall not approve its own decision. Every intent affecting a shared
or contended resource shall be mediated by the conflict fabric before it
can reach an ACT. This invariant is not declarable, admits no exception,
and applies to every DE including those admitted as extensions.
Contention shall be determined by the platform from the action spaces
declared by admitted agents and shall never be asserted by the emitting
agent; where an intent touches only an action space held exclusively by
one agent, the fabric shall return allow without an arbitration cycle.}
\reqfield{Dependency}{%
Req-T23-DE-01}
\reqfield{Traceability (Forward)}{%
\begin{itemize}
\tightlist
\item
  Structural: the DE runtime has no path to ACT except through the
  fabric binding.
\end{itemize}}
\reqfield{UNITY-6G components}{%
\begin{itemize}
\tightlist
\item
  DE runtime
\item
  Conflict fabric
\end{itemize}}
\reqfield{Linked KPIs}{%
\begin{itemize}
\tightlist
\item
  KPI\#21: Reduce the performance penalty due to conflicts by 20\% by
  means of DLT, hierarchical algorithms and policy-based conflict
  resolution schemes.
\item
  KPI\#26: Increase the ratio of service management and LCM tasks
  resolved by local AE/DE components.
\end{itemize}}
\end{requirement}

\begin{requirement}{Req-T23-DE-03}{MUST}{Mediation as an egress property}
\reqfield{Description / rationale}{%
Intent shall leave the DE through the conflict fabric binding, which
returns allow, modify, defer or deny. The binding shall be established
at activation and shall be a precondition of decision emission; an agent
whose conflict binding is unresolved shall never be granted emission.}
\reqfield{Dependency}{%
Req-T23-DE-02}
\reqfield{Traceability (Forward)}{%
\begin{itemize}
\tightlist
\item
  Binding established at activation by the DMO.
\item
  Emission gated by the platform runtime.
\end{itemize}}
\reqfield{UNITY-6G components}{%
\begin{itemize}
\tightlist
\item
  DE runtime
\item
  DMO
\item
  Conflict fabric
\end{itemize}}
\reqfield{Linked KPIs}{%
\begin{itemize}
\tightlist
\item
  KPI\#21: Reduce the performance penalty due to conflicts by 20\% by
  means of DLT, hierarchical algorithms and policy-based conflict
  resolution schemes.
\item
  KPI\#26: Increase the ratio of service management and LCM tasks
  resolved by local AE/DE components.
\item
  KPI\#4: Real-time reconfiguration of IEEE 802.11 network domain in
  range smaller than 10 ms from taking the decision.
\item
  KPI\#5: Real-time reconfiguration of cellular network domain in the
  range smaller than 10 ms from taking the decision.
\end{itemize}}
\end{requirement}

\begin{requirement}{Req-T23-DE-04}{MUST}{Bounded modification}
\reqfield{Description / rationale}{%
Where the fabric returns modify, the modified intent shall be re-checked
against the agent\textquotesingle s own declared action space and
constraints before egress. A modification falling outside them shall be
treated as a denial and shall trigger the declared unresolved-conflict
behavior.}
\reqfield{Dependency}{%
Req-T23-DE-03}
\reqfield{Traceability (Forward)}{%
\begin{itemize}
\tightlist
\item
  Re-check performed by the DE runtime after mediation.
\end{itemize}}
\reqfield{UNITY-6G components}{%
\begin{itemize}
\tightlist
\item
  DE runtime
\end{itemize}}
\reqfield{Linked KPIs}{%
\begin{itemize}
\tightlist
\item
  KPI\#21: Reduce the performance penalty due to conflicts by 20\% by
  means of DLT, hierarchical algorithms and policy-based conflict
  resolution schemes.
\item
  KPI\#26: Increase the ratio of service management and LCM tasks
  resolved by local AE/DE components.
\end{itemize}}
\end{requirement}

\begin{requirement}{Req-T23-DE-05}{MUST}{External enforcement of stability controls}
\reqfield{Description / rationale}{%
Declared stability controls (cooldown, hysteresis, maximum rate of
change, staging) shall be enforced by the platform at the egress gate
and on the actuation path. Intents exceeding a declared control shall be
suppressed and a conformance fault recorded against the agent.}
\reqfield{Dependency}{%
Req-T23-DE-03}
\reqfield{Traceability (Forward)}{%
\begin{itemize}
\tightlist
\item
  Enforced at the egress gate and again on the actuation path.
\end{itemize}}
\reqfield{UNITY-6G components}{%
\begin{itemize}
\tightlist
\item
  DE runtime
\item
  ACT path
\end{itemize}}
\reqfield{Linked KPIs}{%
\begin{itemize}
\tightlist
\item
  KPI\#21: Reduce the performance penalty due to conflicts by 20\% by
  means of DLT, hierarchical algorithms and policy-based conflict
  resolution schemes.
\item
  KPI\#36: Guarantee 99.99999\% service availability and reliability
  based on effective monitoring and LCM.
\item
  KPI\#4: Real-time reconfiguration of IEEE 802.11 network domain in
  range smaller than 10 ms from taking the decision.
\item
  KPI\#5: Real-time reconfiguration of cellular network domain in the
  range smaller than 10 ms from taking the decision.
\end{itemize}}
\end{requirement}

\begin{requirement}{Req-T23-DE-06}{MUST}{Evidence gating before decision logic}
\reqfield{Description / rationale}{%
The DE runtime shall evaluate input quality before the contributed
decision logic runs. Where the degraded-input flag is set or coverage
falls below the declared minimum, the declared degraded-evidence
behavior shall execute and the contributed logic shall not be invoked
for that cycle.}
\reqfield{Dependency}{%
Req-T23-AE-03}
\reqfield{Traceability (Forward)}{%
\begin{itemize}
\tightlist
\item
  Gate precedes contributed logic; not skippable by the agent.
\end{itemize}}
\reqfield{UNITY-6G components}{%
\begin{itemize}
\tightlist
\item
  DE runtime
\item
  AE runtime
\end{itemize}}
\reqfield{Linked KPIs}{%
\begin{itemize}
\tightlist
\item
  KPI\#42: DT-validated AI algorithms that do not break the
  time-sensitive operation when deployed in 99\% of the outputs.
\item
  KPI\#36: Guarantee 99.99999\% service availability and reliability
  based on effective monitoring and LCM.
\item
  KPI\#37: Guarantee 99.99999\% service continuity through AI-driven
  proactive management of cloud-native services.
\end{itemize}}
\end{requirement}

\begin{requirement}{Req-T23-DE-07}{MUST}{Mandatory structured rationale}
\reqfield{Description / rationale}{%
Every emitted intent shall carry a rationale containing the identifiers
of the insight envelopes consumed, the constraints applied, the
objectives targeted and a reason code. Free text alone shall not be
conformant and an intent without a rationale shall be rejected at the
publication boundary. The descriptor shall declare whether the decision
logic is interpretable, meaning that it can produce, for an individual
decision, an account of which inputs determined the outcome and in what
direction; an agent declaring it shall supply that account on request,
and an instance whose descriptor does not declare it shall have an
authority ceiling no higher than
closed\_\allowbreak{}loop\_\allowbreak{}reduced.\allowbreak{}}
\reqfield{Dependency}{%
Req-T23-DE-01}
\reqfield{Traceability (Forward)}{%
\begin{itemize}
\tightlist
\item
  Validated at the publication boundary against the envelope schema.
\end{itemize}}
\reqfield{UNITY-6G components}{%
\begin{itemize}
\tightlist
\item
  DE runtime
\item
  Tuple registry
\end{itemize}}
\reqfield{Linked KPIs}{%
\begin{itemize}
\tightlist
\item
  KPI\#42: DT-validated AI algorithms that do not break the
  time-sensitive operation when deployed in 99\% of the outputs.
\item
  KPI\#21: Reduce the performance penalty due to conflicts by 20\% by
  means of DLT, hierarchical algorithms and policy-based conflict
  resolution schemes.
\end{itemize}}
\end{requirement}

\begin{requirement}{Req-T23-DE-08}{MUST}{Bounded re-decision on actuation feedback}
\reqfield{Description / rationale}{%
ACT outcomes shall be consumed by the DE as evidence for re-decision.
The failure policy shall declare a maximum re-decision count per
objective per window, the counter shall be held by the platform runtime
rather than the agent, and on exhaustion the DE shall stop and raise a
lifecycle event to the DMO recording the objective as unachievable.}
\reqfield{Dependency}{%
Req-T23-DE-06}
\reqfield{Traceability (Forward)}{%
\begin{itemize}
\tightlist
\item
  Counter held by the platform runtime.
\item
  Lifecycle event raised to the DMO on exhaustion.
\end{itemize}}
\reqfield{UNITY-6G components}{%
\begin{itemize}
\tightlist
\item
  DE runtime
\item
  DMO
\item
  ACT runtime
\end{itemize}}
\reqfield{Linked KPIs}{%
\begin{itemize}
\tightlist
\item
  KPI\#31: Reduce service recovery time below 180 s.
\item
  KPI\#36: Guarantee 99.99999\% service availability and reliability
  based on effective monitoring and LCM.
\end{itemize}}
\end{requirement}

\begin{requirement}{Req-T23-DE-09}{MUST}{Declared failure behavior and silence as default}
\reqfield{Description / rationale}{%
The DE failure policy shall declare responses to low confidence with its
threshold, to degraded evidence independently of that threshold, to
unresolved conflict on arbitration timeout, to actuation failure
including the re-decision maximum, and to unavailable insight. Where a
declared behavior is ambiguous or absent at runtime the DE shall emit no
intent; no action shall always be a permitted outcome.}
\reqfield{Dependency}{%
Req-T23-DE-06}
\reqfield{Traceability (Forward)}{%
\begin{itemize}
\tightlist
\item
  Checked by the validity rules.
\item
  Enforced by the DE runtime.
\end{itemize}}
\reqfield{UNITY-6G components}{%
\begin{itemize}
\tightlist
\item
  Tuple descriptor
\item
  DE runtime
\end{itemize}}
\reqfield{Linked KPIs}{%
\begin{itemize}
\tightlist
\item
  KPI\#31: Reduce service recovery time below 180 s.
\item
  KPI\#36: Guarantee 99.99999\% service availability and reliability
  based on effective monitoring and LCM.
\item
  KPI\#37: Guarantee 99.99999\% service continuity through AI-driven
  proactive management of cloud-native services.
\end{itemize}}
\end{requirement}

\begin{requirement}{Req-T23-DE-10}{MUST}{Per-domain decision topology}
\reqfield{Description / rationale}{%
DEs shall be per domain and no runtime global DE shall exist. A
cross-domain agent shall be a federation of per-domain loops, each
bounded by the resource-facing service its domain received. Coherence
shall come from one-time decomposition at order time and from runtime
arbitration by the conflict fabric.}
\reqfield{Dependency}{%
Req-T23-DE-02}
\reqfield{Traceability (Forward)}{%
\begin{itemize}
\tightlist
\item
  Structural; enforced at admission by domain assignment.
\end{itemize}}
\reqfield{UNITY-6G components}{%
\begin{itemize}
\tightlist
\item
  Tuple registry
\item
  IDMO
\item
  DMO
\end{itemize}}
\reqfield{Linked KPIs}{%
\begin{itemize}
\tightlist
\item
  KPI\#26: Increase the ratio of service management and LCM tasks
  resolved by local AE/DE components.
\item
  KPI\#21: Reduce the performance penalty due to conflicts by 20\% by
  means of DLT, hierarchical algorithms and policy-based conflict
  resolution schemes.
\item
  KPI\#1: Deliver four distributed management entities (MS, AE, DE and
  ACT) deployable in each heterogeneous domain.
\end{itemize}}
\end{requirement}

\begin{requirement}{Req-T23-DE-11}{MUST}{Lateral decisions do not bind}
\reqfield{Description / rationale}{%
A decision exchanged between domain agents shall be a proposal
conferring no authority over the receiving agent, which shall remain
answerable for what it applies within its own scope. A binding
cross-domain decision shall be taken at the cross-domain management
level. Authority shall descend and shall not pass laterally.}
\reqfield{Dependency}{%
Req-T23-DE-10}
\reqfield{Traceability (Forward)}{%
\begin{itemize}
\tightlist
\item
  Enforced by the egress gate and the action space of the receiving
  agent.
\end{itemize}}
\reqfield{UNITY-6G components}{%
\begin{itemize}
\tightlist
\item
  DE runtime
\item
  IDMO
\item
  Conflict fabric
\end{itemize}}
\reqfield{Linked KPIs}{%
\begin{itemize}
\tightlist
\item
  KPI\#21: Reduce the performance penalty due to conflicts by 20\% by
  means of DLT, hierarchical algorithms and policy-based conflict
  resolution schemes.
\item
  KPI\#26: Increase the ratio of service management and LCM tasks
  resolved by local AE/DE components.
\end{itemize}}
\end{requirement}

\begin{requirement}{Req-T23-DE-12}{MUST}{Identical containment for extension agents}
\reqfield{Description / rationale}{%
A DE admitted as an extension shall be granted reduced authority through
narrower control scope, tighter action space bounds and lower
arbitration priority, and shall use the identical egress path, conflict
mediation, stability enforcement and rationale requirements as any other
agent. Sustained conformant operation under evaluation shall be the
route to promotion.}
\reqfield{Dependency}{%
Req-T23-DE-02}
\reqfield{Traceability (Forward)}{%
\begin{itemize}
\tightlist
\item
  Admitted at reduced authority; identical guardrails applied by the
  runtime.
\end{itemize}}
\reqfield{UNITY-6G components}{%
\begin{itemize}
\tightlist
\item
  Admission (IDMO/DMO)
\item
  Evaluation function
\item
  DE runtime
\end{itemize}}
\reqfield{Linked KPIs}{%
\begin{itemize}
\tightlist
\item
  KPI\#26: Increase the ratio of service management and LCM tasks
  resolved by local AE/DE components.
\item
  KPI\#41: Reduction of 70\% of deployment time of AI-based
  resource/energy management algorithms.
\end{itemize}}
\end{requirement}

\hypertarget{contribution-to-the-project-kpis-2}{%
\subsection{CONTRIBUTION TO THE PROJECT KPIs}\label{contribution-to-the-project-kpis-2}}

Table 6-4 states what the Decision Engine contributes to each KPI it
bears on, and Table 6-5 maps the requirements of this chapter onto those
KPIs.

\begin{wptable}{6-4}

\begin{longtable}[]{@{}
  >{\raggedright\arraybackslash}p{(\columnwidth - 4\tabcolsep) * \real{0.1499}}
  >{\raggedright\arraybackslash}p{(\columnwidth - 4\tabcolsep) * \real{0.6375}}
  >{\raggedright\arraybackslash}p{(\columnwidth - 4\tabcolsep) * \real{0.2126}}@{}}
\caption{Contribution of the Decision Engine to the project
KPIs}\tabularnewline
\toprule\noalign{}
\begin{minipage}[b]{\linewidth}\raggedright
\textbf{KPI}
\end{minipage} & \begin{minipage}[b]{\linewidth}\raggedright
\textbf{Contribution of the Decision Engine}
\end{minipage} & \begin{minipage}[b]{\linewidth}\raggedright
\textbf{Requirements}
\end{minipage} \\
\midrule\noalign{}
\endfirsthead
\toprule\noalign{}
\begin{minipage}[b]{\linewidth}\raggedright
\textbf{KPI}
\end{minipage} & \begin{minipage}[b]{\linewidth}\raggedright
\textbf{Contribution of the Decision Engine}
\end{minipage} & \begin{minipage}[b]{\linewidth}\raggedright
\textbf{Requirements}
\end{minipage} \\
\midrule\noalign{}
\endhead
\bottomrule\noalign{}
\endlastfoot
\textbf{KPI\#1} & Supplies the DE of the four-entity tuple, per domain,
with authority bounded by the platform rather than by contributed logic.
& Req-T23-DE-01, Req-T23-DE-10 \\
\textbf{KPI\#4} & The egress gate and stability enforcement sit on the
fast path, so sub-10 ms reconfiguration remains mediated rather than
unmediated. & Req-T23-DE-03, Req-T23-DE-05 \\
\textbf{KPI\#5} & The same gate governs cellular reconfiguration. &
Req-T23-DE-03, Req-T23-DE-05 \\
\textbf{KPI\#21} & Conflict-driven performance penalty is reduced
because no decision reaches actuation unmediated and modification is
bounded by the agent\textquotesingle s own envelope. & Req-T23-DE-02,
Req-T23-DE-03, Req-T23-DE-04 \\
\textbf{KPI\#42} & Every intent carries structured rationale naming the
evidence consumed, the constraints applied and a reason code, and
decision logic never runs on evidence below its declared quality, so an
AI-formed decision can be validated and attributed before it touches
time-sensitive operation. & Req-T23-DE-01, Req-T23-DE-06,
Req-T23-DE-07 \\
\textbf{KPI\#26} & Per-domain decision topology is what allows a task to
be resolved by local components rather than escalated. & Req-T23-DE-10,
Req-T23-DE-11 \\
\textbf{KPI\#31} & Bounded re-decision converts an unachievable
objective into a lifecycle event within the recovery budget rather than
an indefinite loop. & Req-T23-DE-08, Req-T23-DE-09 \\
\textbf{KPI\#36} & Silence as the default on ambiguous failure behavior,
and externally enforced stability controls, protect availability against
oscillating agents. & Req-T23-DE-05, Req-T23-DE-09 \\
\textbf{KPI\#37} & Degraded evidence is gated before decision logic
runs, so proactive action is never taken on evidence that did not
support it. & Req-T23-DE-06 \\
\textbf{KPI\#41} & Extension agents are admitted without bespoke
integration and at identical containment, shortening deployment of new
algorithms. & Req-T23-DE-12 \\
\end{longtable}

\end{wptable}

\begin{wptable}{6-5}\footnotesize

\begin{longtable}[]{@{}
  >{\raggedright\arraybackslash}p{(\columnwidth - 20\tabcolsep) * \real{0.1814}}
  >{\raggedright\arraybackslash}p{(\columnwidth - 20\tabcolsep) * \real{0.0819}}
  >{\raggedright\arraybackslash}p{(\columnwidth - 20\tabcolsep) * \real{0.0819}}
  >{\raggedright\arraybackslash}p{(\columnwidth - 20\tabcolsep) * \real{0.0819}}
  >{\raggedright\arraybackslash}p{(\columnwidth - 20\tabcolsep) * \real{0.0819}}
  >{\raggedright\arraybackslash}p{(\columnwidth - 20\tabcolsep) * \real{0.0819}}
  >{\raggedright\arraybackslash}p{(\columnwidth - 20\tabcolsep) * \real{0.0819}}
  >{\raggedright\arraybackslash}p{(\columnwidth - 20\tabcolsep) * \real{0.0819}}
  >{\raggedright\arraybackslash}p{(\columnwidth - 20\tabcolsep) * \real{0.0819}}
  >{\raggedright\arraybackslash}p{(\columnwidth - 20\tabcolsep) * \real{0.0819}}
  >{\raggedright\arraybackslash}p{(\columnwidth - 20\tabcolsep) * \real{0.0814}}@{}}
\caption{Mapping of Decision Engine requirements onto the project
KPIs}\tabularnewline
\toprule\noalign{}
\begin{minipage}[b]{\linewidth}\raggedright
\textbf{KPI\textbackslash{}}

\textbf{Requirement}
\end{minipage} & \begin{minipage}[b]{\linewidth}\raggedright
\textbf{\#1}
\end{minipage} & \begin{minipage}[b]{\linewidth}\raggedright
\textbf{\#4}
\end{minipage} & \begin{minipage}[b]{\linewidth}\raggedright
\textbf{\#5}
\end{minipage} & \begin{minipage}[b]{\linewidth}\raggedright
\textbf{\#21}
\end{minipage} & \begin{minipage}[b]{\linewidth}\raggedright
\textbf{\#26}
\end{minipage} & \begin{minipage}[b]{\linewidth}\raggedright
\textbf{\#31}
\end{minipage} & \begin{minipage}[b]{\linewidth}\raggedright
\textbf{\#36}
\end{minipage} & \begin{minipage}[b]{\linewidth}\raggedright
\textbf{\#37}
\end{minipage} & \begin{minipage}[b]{\linewidth}\raggedright
\textbf{\#41}
\end{minipage} & \begin{minipage}[b]{\linewidth}\raggedright
\textbf{\#42}
\end{minipage} \\
\midrule\noalign{}
\endfirsthead
\toprule\noalign{}
\begin{minipage}[b]{\linewidth}\raggedright
\textbf{KPI\textbackslash{}}

\textbf{Requirement}
\end{minipage} & \begin{minipage}[b]{\linewidth}\raggedright
\textbf{\#1}
\end{minipage} & \begin{minipage}[b]{\linewidth}\raggedright
\textbf{\#4}
\end{minipage} & \begin{minipage}[b]{\linewidth}\raggedright
\textbf{\#5}
\end{minipage} & \begin{minipage}[b]{\linewidth}\raggedright
\textbf{\#21}
\end{minipage} & \begin{minipage}[b]{\linewidth}\raggedright
\textbf{\#26}
\end{minipage} & \begin{minipage}[b]{\linewidth}\raggedright
\textbf{\#31}
\end{minipage} & \begin{minipage}[b]{\linewidth}\raggedright
\textbf{\#36}
\end{minipage} & \begin{minipage}[b]{\linewidth}\raggedright
\textbf{\#37}
\end{minipage} & \begin{minipage}[b]{\linewidth}\raggedright
\textbf{\#41}
\end{minipage} & \begin{minipage}[b]{\linewidth}\raggedright
\textbf{\#42}
\end{minipage} \\
\midrule\noalign{}
\endhead
\bottomrule\noalign{}
\endlastfoot
\textbf{Req-T23-DE-01} & \kpimark{} & & & & & & & & & \kpimark{} \\
\textbf{Req-T23-DE-02} & & & & \kpimark{} & \kpimark{} & & & & & \\
\textbf{Req-T23-DE-03} & & \kpimark{} & \kpimark{} & \kpimark{} &
\kpimark{} & & & & & \\
\textbf{Req-T23-DE-04} & & & & \kpimark{} & \kpimark{} & & & & & \\
\textbf{Req-T23-DE-05} & & \kpimark{} & \kpimark{} & \kpimark{} & & &
\kpimark{} & & & \\
\textbf{Req-T23-DE-06} & & & & & & & \kpimark{} & \kpimark{} & &
\kpimark{} \\
\textbf{Req-T23-DE-07} & & & & \kpimark{} & & & & & & \kpimark{} \\
\textbf{Req-T23-DE-08} & & & & & & \kpimark{} & \kpimark{} & & & \\
\textbf{Req-T23-DE-09} & & & & & & \kpimark{} & \kpimark{} & \kpimark{}
& & \\
\textbf{Req-T23-DE-10} & \kpimark{} & & & \kpimark{} & \kpimark{} & & &
& & \\
\textbf{Req-T23-DE-11} & & & & \kpimark{} & \kpimark{} & & & & & \\
\textbf{Req-T23-DE-12} & & & & & \kpimark{} & & & & \kpimark{} & \\
\end{longtable}

\end{wptable}
